\chapter{トレンドフォロー}
\label{ch:trend}

\section{戦略の中身}

トレンドフォロー（時系列モメンタム、マネージド・フューチャーズ、CTA の戦略）は、株価指数・国債・通貨・商品の先物を、それ自身の過去のリターンが正なら買い、負なら売る。

\begin{equation}
w_{i,t} = \mathrm{sign}\Big(\sum_{k=1}^{12} r_{i,t-k}\Big) \times \frac{\sigma^*}{\hat\sigma_{i,t}}
\end{equation}

\begin{itemize}
\item \textbf{信号}：過去 1〜12 か月のリターンの符号、移動平均の交差（短期の平均が長期の平均を上回る）など。期間の違う信号を組み合わせることが多い。
\item \textbf{量}：各市場の変動率 $\hat\sigma_{i,t}$ の逆数に比例させ、どの市場も同じリスクにする（ボラティリティ・スケーリング）。
\item \textbf{分散}：数十〜百以上の先物市場に分散する。
\end{itemize}

横断面のモメンタム（第 \ref{ch:equityfactor} 章、銘柄を比べて勝ち組を買う）と違い、時系列のモメンタムは、すべての市場が上がっていればすべて買う。

\section{中心の論文}

\begin{paperbox}[Moskowitz・Ooi・Pedersen（2012）「時系列のモメンタム」\cite{mop2012}]
58 の流動性の高い先物・フォワード（株価指数、通貨、商品、国債）で、各資産の過去 12 か月のリターンが、その後 1 か月のリターンを正に予測することを示した。予測の力は 1〜12 か月で続き、より長い期間では一部が戻る。全資産に分散した時系列モメンタムの戦略は、市場・割安・横断面のモメンタムなどの標準的なファクターでは説明できない大きな超過リターンを上げ、市場が極端に上げ下げした時期に最も良い成績を上げた。投機筋（非商業の先物の参加者）がトレンドに乗り、ヘッジをする当業者がその反対側に立っていることも示した。
\end{paperbox}

\begin{paperbox}[Hurst・Ooi・Pedersen（2017）「トレンドフォローの 1 世紀の証拠」\cite{hop2017}]
1880 年まで遡って、株価指数・国債・通貨・商品の時系列モメンタムの戦略を作り直し、ほぼすべての 10 年間で正の平均リターンを上げたこと、景気後退や戦争・金融危機を含む多くの環境で成績を保ったこと、株式や債券の大きな下げの時期に利益を上げることが多かったことを示した。
\end{paperbox}

\section{なぜ儲かるのか}

\begin{itemize}
\item \textbf{反応の不足と、その後の行き過ぎ}：新しい情報への反応が遅れ（最初の反応の不足）、その後、群集の行動で行き過ぎる。どちらの局面でも、トレンドに乗れば利益が出る。
\item \textbf{当業者のヘッジ}：生産者・消費者が価格のリスクを避けるために先物を売買し、その反対に立つ投機筋に対価が支払われる。
\item \textbf{中央銀行の介入}：為替や金利の政策は、ゆっくり一方向に動くことが多い。
\end{itemize}

\section{オプションのような形}

Fung・Hsieh（2001）\cite{funghsieh2001} は、トレンドフォローの運用者のリターンが、「期間の中の最高値で売り、最安値で買う」権利を持つオプション（ルックバック・ストラドル）を買う戦略のリターンに似ていることを示した。トレンドフォローは、大きく一方向に動いた市場で利益を上げ、方向のない市場では小さな損を重ねる。オプションを買い続けているのと同じ形で、リターンの分布は正の歪度を持つ。この「ルックバック・ストラドルのファクター」は、ヘッジファンドの成績の評価に使われる Fung・Hsieh の 7 ファクター・モデル（第 \ref{ch:hfperf} 章）の一部になった。

\section{批判}

\begin{itemize}
\item \textbf{ボラティリティ・スケーリングの寄与}：Kim・Tse・Wald（2016）\cite{kimtsewald2016} は、時系列モメンタムの成績の大部分は、各市場の量を変動率で調整すること（ボラティリティ・スケーリング）から来ていて、量を調整しなければ、単に買い持ちする戦略と大差ないと論じた。
\item \textbf{予測の力は本当にあるか}：Huang・Li・Wang・Zhou（2020）\cite{hlwz2020} は、資産ごとに見ると、過去 12 か月のリターンが将来のリターンを予測する証拠は弱く、全資産をまとめた回帰の有意性は一部の資産と推定の偏りに依存すると論じた。時系列モメンタムの戦略の収益は、各資産の長期の平均リターン（予測の力がなくても得られる部分）でかなり説明できる、という。
\end{itemize}

\section{弱点}

\begin{itemize}
\item \textbf{方向のない相場}：上げ下げを繰り返す相場では、信号が入れ替わるたびに損をする（往復の損）。
\item \textbf{急な反転}：トレンドが急に反転すると、反応が遅れて損をする。2009 年の春、2020 年 3 月の急落と急反発など。
\item \textbf{長い低迷}：2010 年代は、金利が低く変動の小さい環境で、多くの CTA の成績が低迷した。物価と金利が大きく動いた 2022 年には、多くが大きな利益を上げた（業界の指数で広く報じられた事実で、特定の論文の値ではない）。
\end{itemize}

\section{個人や研究で再現できるか}

\begin{itemize}
\item 先物のつなぎ足の価格のデータがあれば、検証は難しくない。規則が単純で、過学習の余地も比較的小さい。
\item 実運用では、数十の先物市場に分散するための証拠金、限月の乗り換え（ロール）のコスト、最小の取引単位が壁になる。個人なら、先物の代わりに ETF で近似することになる。
\item 結果は、信号の期間・ボラティリティの推定の窓・市場の選び方に敏感である。多くの組み合わせを試すと、過学習（\chapref{ch:backtest}）の危険がある。
\end{itemize}

\begin{projbox}
トレンドフォローは、LLM-BLM の研究とは別の種類の情報（テキストではなく価格）を使うが、「予測の力そのものより、量の調整（ボラティリティ・スケーリング）とリスクの配分で成績の多くが決まる」という論争（Kim ほか 2016、Huang ほか 2020）は、LLM の点数を使う運用にもそのまま当てはまる。点数の予測の力を評価するときに、ポートフォリオの作り方の効果と分けて測る必要がある。
\end{projbox}
